%%%%%%%%%%%%%%%%%%%%%%%%%%%%%%%%%%%%%%%%%
% Jacobs Landscape Poster
% LaTeX Template
% Version 1.0 (29/03/13)
%
% Created by:
% Computational Physics and Biophysics Group, Jacobs University
% https://teamwork.jacobs-university.de:8443/confluence/display/CoPandBiG/LaTeX+Poster
%
% Further modified by:
% Nathaniel Johnston (nathaniel@njohnston.ca)
%
% This template has been downloaded from:
% http://www.LaTeXTemplates.com
%
% License:
% CC BY-NC-SA 3.0 (http://creativecommons.org/licenses/by-nc-sa/3.0/)
%
%%%%%%%%%%%%%%%%%%%%%%%%%%%%%%%%%%%%%%%%%

%----------------------------------------------------------------------------------------
%	PACKAGES AND OTHER DOCUMENT CONFIGURATIONS
%----------------------------------------------------------------------------------------

\documentclass[final]{beamer}

\usepackage[scale=1.24]{beamerposter} % Use the beamerposter package for laying out the poster
\usepackage{float}
\usepackage{onimage}

\usetheme{confposter} % Use the confposter theme supplied with this template

\setbeamercolor{block title}{fg=ngreen,bg=white} % Colors of the block titles
\setbeamercolor{block body}{fg=black,bg=white} % Colors of the body of blocks
\setbeamercolor{block alerted title}{fg=white,bg=dblue!70} % Colors of the highlighted block titles
\setbeamercolor{block alerted body}{fg=black,bg=dblue!10} % Colors of the body of highlighted blocks
% Many more colors are available for use in beamerthemeconfposter.sty

%-----------------------------------------------------------
% Define the column widths and overall poster size
% To set effective sepwid, onecolwid and twocolwid values, first choose how many columns you want and how much separation you want between columns
% In this template, the separation width chosen is 0.024 of the paper width and a 4-column layout
% onecolwid should therefore be (1-(# of columns+1)*sepwid)/# of columns e.g. (1-(4+1)*0.024)/4 = 0.22
% Set twocolwid to be (2*onecolwid)+sepwid = 0.464
% Set threecolwid to be (3*onecolwid)+2*sepwid = 0.708

\newlength{\sepwid}
\newlength{\onecolwid}
\newlength{\twocolwid}
\newlength{\threecolwid}
\setlength{\paperwidth}{48in} % A0 width: 46.8in
\setlength{\paperheight}{36in} % A0 height: 33.1in
\setlength{\sepwid}{0.024\paperwidth} % Separation width (white space) between columns
\setlength{\onecolwid}{0.22\paperwidth} % Width of one column
\setlength{\twocolwid}{0.464\paperwidth} % Width of two columns
\setlength{\threecolwid}{0.708\paperwidth} % Width of three columns
\setlength{\topmargin}{-0.5in} % Reduce the top margin size
%-----------------------------------------------------------

\usepackage{graphicx}  % Required for including images

\usepackage{booktabs} % Top and bottom rules for tables

% Macros for arrows between images
% based on http://tex.stackexchange.com/questions/7925/big-arrows-between-images
%----------------------------------------------------------------------------------------
\newcommand*{\vimage}[1]{\vcenter{\hbox{#1}}}
\newcommand*{\vpointer}[1][\ \triangleright\ ]{\vcenter{\hbox{$#1$}}}
\usepackage{mathtools} % for xRightarrow

% Tikz stuff for the cool diagrams
%----------------------------------------------------------------------------------------
\usetikzlibrary{spy}

%----------------------------------------------------------------------------------------
%	TITLE SECTION
%----------------------------------------------------------------------------------------

\title{Finding Eigenvalues using Automatic Differentiation} % Poster title

\author{Alex Hirzel, Peter Solfest, and Ryan Bruner} % Author(s)

\institute{Michigan Technological University} % Institution(s)

%----------------------------------------------------------------------------------------

\begin{document}

\addtobeamertemplate{block end}{}{\vspace*{2ex}} % White space under blocks
\addtobeamertemplate{block alerted end}{}{\vspace*{2ex}} % White space under highlighted (alert) blocks

\setlength{\belowcaptionskip}{2ex} % White space under figures
\setlength\belowdisplayshortskip{2ex} % White space under equations

\begin{frame}[t] % The whole poster is enclosed in one beamer frame

\begin{columns}[t] % The whole poster consists of three major columns, the second of which is split into two columns twice - the [t] option aligns each column's content to the top

\begin{column}{\sepwid}\end{column} % Empty spacer column

\begin{column}{\onecolwid} % The first column

%----------------------------------------------------------------------------------------
%	ABSTRACT
%----------------------------------------------------------------------------------------

\begin{block}{Abstract}
{
\small
Eigenvalues are ubiquitous in all branches of science and engineering. The
dominant eigenvalue algorithm reduces input matrices to almost diagonal and then
creates and chases bulges using parameters called shifts. Good parameter choices
can allow shrinking of the original problem allowing for a more aggressive
deflation on the subproblem. Such deflations greatly improve performance. We
apply automatic differentiation tools to achieve frequent deflations.
}
\end{block}

%----------------------------------------------------------------------------------------
%	EXISTING TECHNIQUES
%----------------------------------------------------------------------------------------

\begin{block}{Existing Techniques}

\vspace{0.5ex}
\begin{figure}
$\vimage{\includegraphics[width=0.2\linewidth]{plots/original.eps}}
\vpointer
\vimage{\includegraphics[width=0.2\linewidth]{plots/qr1.eps}}
\vpointer
\vimage{\includegraphics[width=0.2\linewidth]{plots/qr2.eps}}
\vpointer
\vimage{\includegraphics[width=0.2\linewidth]{plots/qrdone.eps}}$
\end{figure}
\vspace{0.5ex}

A general matrix is first reduced to Hessenberg~form, then the $\mathbf{Q}
\mathbf{R}$~\textbf{algorithm} is used to reduce the sub-diagonal entries to
zero (shown above). This exposes eigenvalues on the diagonal. The $Q
R$~algorithm consists of repeated iterations of:
%
\begin{align}
A_n     &\rightarrow Q_n \cdot R_n \label{eq:qr} \\
A_{n+1} &\leftarrow  R_n \cdot Q_n = Q_n^{-1} \cdot A_n \cdot Q_n \label{eq:rq}
\end{align}
%
As written, \eqref{eq:qr} and \eqref{eq:rq} require matrix multiplies which can
be very large. For this reason, the $Q R$~algorithm is implemented
\emph{implicitly} as a loop of similarity transformations. This results in a
\emph{bulge} moving from top-left to bottom-right:

\vspace{0.5ex}
\begin{figure}
$\vimage{\includegraphics[width=0.2\linewidth]{plots/bulge1.eps}}
\vpointer
\vimage{\includegraphics[width=0.2\linewidth]{plots/bulge2.eps}}
\vpointer
\vimage{\includegraphics[width=0.2\linewidth]{plots/bulge3.eps}}
\vpointer
\vimage{\includegraphics[width=0.2\linewidth]{plots/bulgedone.eps}}$
\end{figure}
\vspace{0.5ex}

This bulge is seeded with a starting value, a \emph{shift}. A second bulge can
be started from the bottom right. When bulges are chased from each side, they
meet in the middle. A Schur decomposition is used to form \emph{spikes} around
this bulge. When two or more of these \emph{spike values} are sufficiently
close to zero, the problem can be subdivided recursively into the top-left
and bottom-right halves. This is shown below.

\vspace{0.5ex}
\begin{figure}
$\vimage{\includegraphics[width=0.2\linewidth]{plots/bulges1.eps}}
\vpointer
\vimage{\includegraphics[width=0.2\linewidth]{plots/bulges2.eps}}
\vpointer
\vimage{\includegraphics[width=0.2\linewidth]{plots/bulgesspiked.eps}}
\vpointer
\vimage{\includegraphics[width=0.2\linewidth]{plots/bulgesdeflated.eps}}
$
\end{figure}
\vspace{0.5ex}

\end{block}
\end{column} % End of the first column

\begin{column}{\sepwid}\end{column} % Empty spacer column

\begin{column}{\twocolwid} % Begin a column which is two columns wide (column 2)

\begin{alertblock}{Proposed Technique}

\begin{figure}[H]
$\vimage{\includegraphics[width=0.2\linewidth]{plots2/impstep001.eps}}
\vpointer
\vimage{\includegraphics[width=0.2\linewidth]{plots2/impstep006.eps}}
\vpointer
\vimage{\includegraphics[width=0.2\linewidth]{plots2/impstep008.eps}}
\vpointer
\vimage{\includegraphics[width=0.2\linewidth]{plots2/impstep009.eps}}
$
\end{figure}

Other text TBD

Note that this whole poster is now in English (no Latin :) ).
\end{alertblock}

\begin{columns}[t,totalwidth=\twocolwid] % Split up the two columns wide column

\begin{column}{\onecolwid}\vspace{-.6in} % The first column within column 2 (column 2.1)

%----------------------------------------------------------------------------------------
%	Schur Decomposition?
%----------------------------------------------------------------------------------------
\begin{block}{Schur Decomposition?}
The Schur decomposition is what allows us to create spikes. For an input matrix
$A$, it forms $Q$ and $T$ such that $Q$ diagonalizes $T$ into $A$. In symbols,
%
\begin{align}
A = Q^\intercal \cdot T \cdot Q
\end{align}
%
Since $Q$ diagonalizes $A$, we can consider a matrix $A_\text{sub}$ and compute
$Q_\text{sub}$:
%
\begin{figure}[H]
\begin{align*}
\vimage{ \begin{tikzonimage}[width=0.28\linewidth]{plots/schur_Az.eps}[white]
\node[fill=black, anchor=north east] at (1, 1) {$A_\text{sub}$};
\end{tikzonimage} }
\vpointer
\vimage{ \begin{tikzonimage}[width=0.28\linewidth]{plots/schur_Qz.eps}[white]
\node[fill=black, anchor=north east] at (1, 1) {$Q_\text{sub}$};
\end{tikzonimage} }
\end{align*}
\end{figure}
%
We can then embed $Q_\text{sub}$ into $Q$ and apply it to $A$. In the
neighborhood of the Schur decomposition, the subdiagonal of $A$ will be zero
and the spikes will form.
%
\begin{figure}[H]
\begin{align*}
\vimage{ \begin{tikzonimage}[width=0.28\linewidth]{plots/schur_Q.eps}[white]
\node[fill=black, anchor=north east] at (1, 1) {$Q$};
\end{tikzonimage} }
{\huge \,\times\, }
\vimage{ \begin{tikzonimage}[width=0.28\linewidth]{plots/schur_Ahess.eps}[white]
\node[fill=black, anchor=north east] at (1, 1) {$A$};
\end{tikzonimage} }
&
\vpointer
\vimage{ \begin{tikzonimage}[width=0.28\linewidth]{plots/schur_final.eps}[white]
\node[fill=black, anchor=north east] at (1, 1) {$Q^\intercal \!\cdot\! A \!\cdot\! Q$};
\end{tikzonimage} }
\end{align*}
\end{figure}
%
\end{block}


%----------------------------------------------------------------------------------------

\end{column} % End of column 2.1

\begin{column}{\onecolwid}\vspace{-.6in} % The second column within column 2 (column 2.2)


%----------------------------------------------------------------------------------------
%	Automatic Differentiation?
%----------------------------------------------------------------------------------------
\begin{block}{Automatic Differentiation?}
In practice, middle deflations are improbable because they rely on the values
of the spikes. The spike values depend on the shifts fed to the $Q R$ algorithm. 
are generated using a Schur decomposition.
They occur when the spike tips become zero due to change or luck. Inducing shifts

\vspace{0.5ex}
\begin{figure}
\begin{tikzpicture}[spy using outlines={circle, magnification=3, connect spies}]
\node {\includegraphics[width=0.2\linewidth]{plots/bulgesspiked.eps}};
\spy [ngreen, size=3.5in] on (-0.7,-0.7) in node [left] at (-2in,0);
\end{tikzpicture}
\end{figure}
\vspace{0.5ex}

We apply \textbf{automatic differentiation} (AD) tools to these spikes. AD tools
allow \emph{algorithms} to be differentiated in the same sense as algebraic
functions. Classical calculus tells us

\begin{equation*}
\frac{d}{dx} \left( \smash{\overbracket{x^2 + 2x + 1}^{f(x)}} \right) =
\overbracket{2x + 2}^{f'(x)}\text{,}
\end{equation*}

which allows us to use Newton's method

\begin{equation}
x_{n+1} = x_n - \frac{f(x)}{f'(x)}
\end{equation}

to find the zeros of $f(x)$. \textbf{The technique for finding zeros of the
spike values is analogus.} Applying AD tools to the Schur decomposition creates
the gradient with respect to the input values (shifts). When the proper shifts
are found to induce a zero on the spike, a deflation has been found.
\end{block}


%----------------------------------------------------------------------------------------
%	FUTURE WORK
%----------------------------------------------------------------------------------------
\begin{block}{Future Work}
In this section, cover:
\begin{itemize}
\item What are we going to do next? Very much depends on results that have not
yet been teased out.
\end{itemize}
\end{block}

%----------------------------------------------------------------------------------------

\end{column} % End of column 2.2

\end{columns} % End of the split of column 2 - any content after this will now take up 2 columns width

%----------------------------------------------------------------------------------------

%\begin{columns}[t,totalwidth=\twocolwid] % Split up the two columns wide column again
%
%\begin{column}{\onecolwid} % The first column within column 2 (column 2.1)
%
%\end{column} % End of column 2.1
%
%\begin{column}{\onecolwid} % The second column within column 2 (column 2.2)
%
%%----------------------------------------------------------------------------------------
%
%\end{column} % End of column 2.2
%
%\end{columns} % End of the split of column 2

%----------------------------------------------------------------------------------------
%	PERFORMANCE IMPROVEMENT
%----------------------------------------------------------------------------------------

\begin{alertblock}{Performance Improvement}

Our algorithm outperforms XXX by XXXX. We show this by using a MATLAB package
which counts FLOPs which shows that our algorithm performs fewer operations when
compared to traditional QR on the same matrix.

\begin{figure}[h]
(pretty pictures here)
\end{figure}

Other text TBD

Note that this whole poster is now in English (no Latin :) ).

There is a very cool-looking package that might work with MATLAB for counting
FLOPS:
\url{http://research.microsoft.com/en-us/um/people/minka/software/lightspeed/}

\end{alertblock}

\end{column} % End of the second column

\begin{column}{\sepwid}\end{column} % Empty spacer column

\begin{column}{\onecolwid} % The third column

%----------------------------------------------------------------------------------------
%	CONCLUSION
%----------------------------------------------------------------------------------------

\begin{block}{Conclusion}

We successfully apply Automatic Differentiation techniques to induce deflations
in the center of matrices. With the proper implementation, this will provide a
large performance improvement over traditional algorithms in use today.
Additionally, this is an enabling technique for extremely large eigenvalue
problems. Our preliminary results indicate that a performance gain of XXX can be
expected from proper implementations of this technique. This was determined by
naively implementing each algorithm in MATLAB and comparing FLOP counts using
each algorithm for the same inputs.

\end{block}

%----------------------------------------------------------------------------------------
%	ADDITIONAL INFORMATION
%----------------------------------------------------------------------------------------

\begin{block}{Additional Information}

The authors are pursuing publication. For more details on this technique, please
contact the authors directly.

\end{block}

%----------------------------------------------------------------------------------------
%	REFERENCES
%----------------------------------------------------------------------------------------

\begin{block}{References}

\nocite{*} % Insert publications even if they are not cited in the poster
\small{\bibliographystyle{siam}
\bibliography{papers}\vspace{0.75in}}

\end{block}

%----------------------------------------------------------------------------------------
%	ACKNOWLEDGEMENTS
%----------------------------------------------------------------------------------------

%\setbeamercolor{block title}{fg=red,bg=white} % Change the block title color

\begin{block}{Acknowledgements}

The authors thank \textbf{Dr.~Allan~Struthers} (Mathematics
department) for his guidance.

\end{block}

%----------------------------------------------------------------------------------------
%	CONTACT INFORMATION
%----------------------------------------------------------------------------------------

\setbeamercolor{block alerted title}{fg=black,bg=norange} % Change the alert block title colors
\setbeamercolor{block alerted body}{fg=black,bg=white} % Change the alert block body colors

\begin{block}{Contact Information}

\begin{itemize}
%\item Web: \href{http://www.university.edu/smithlab}{http://www.university.edu/smithlab}
\item Email: \href{mailto:ahirzel@mtu.edu}{\texttt{ahirzel@mtu.edu}}
\item Phone: +1 (906) 231 0866
\end{itemize}

\end{block}

\begin{center}
% minimum size is 1 3/4"
\includegraphics[width=0.4\linewidth]{logo_withfullname_LG.eps}
\end{center}

%----------------------------------------------------------------------------------------

\end{column} % End of the third column

\end{columns} % End of all the columns in the poster

\end{frame} % End of the enclosing frame

\end{document}
